\documentclass[11pt]{article}
\usepackage[margin=25mm]{geometry}
\usepackage{amsmath,amssymb}
\begin{document}

\title{I01 prevalidation 16 result:\\Leading angle law of paired support transfer}
\date{2026-10-04}
\maketitle

\section*{Leading response}
For the three nonzero $n=3$ classes, the contrast divided by $\sin^2\epsilon$
over eight non-endpoint angles has relative spread $2.61\times10^{-5}$,
$9.41\times10^{-6}$, and $9.52\times10^{-6}$ for same support, different-pair
same-sign, and different-pair opposite-sign transfer. Their leading scales are
approximately $2.0367\times10^{-3}$, $5.4742\times10^{-4}$, and
$9.8550\times10^{-4}$, respectively. The conjugate-partner class stays at the
numerical-zero scale throughout.

\section*{Endpoint qualification}
At $\epsilon=\pi$, a small but backend-independent residual remains:
$9.74\times10^{-8}$ for same support, $9.42\times10^{-9}$ for different-pair
same-sign, and $1.72\times10^{-8}$ for different-pair opposite-sign. An independent
raw Qiskit statevector calculation agrees with Aer for the same-support endpoint to
$7.6\times10^{-15}$ in probabilities, so this is not caused by the standard-gate
decomposition used for Aer execution.

\section*{Interpretation}
The observed law is a dominant $\sin^2\epsilon$ transverse-interference response,
not an exact one-term identity at the tested precision. The endpoint residual must
be retained when making any closed-form claim. This motivates finite-shot testing of
the support classes at $\epsilon=\pi/2$, where the leading contrasts are largest.

\section*{Reproduction}
Code: \texttt{../src/relational\_atlas\_i01/holdout/qiskit\_rank\_holdout16\_angle\_law.py}.\\
Output: \texttt{../results/qiskit\_rank\_holdout16\_angle\_law.json}.
\end{document}
